% Round 15. A complete PWS-represented calculation on projective space.
% Uses constructive PWS three-term boundary certificates, not arbitrary actual trace.
\section{The PWS-represented Aeppli groups of projective space in every bidegree}
\label{r15allps:section}

The symmetric-product finite cover gives a shorter route to the
PWS-represented calculation on projective space in all degrees.
The proof retains a stronger boundary certificate than mere
vanishing in the actual Aeppli quotient: its two first-derivative
primitive inputs and its separate $D$-primitive input are PWS.
This is exactly what permits the established all-bidegree finite
PWS trace to descend the certificate.

\begin{definition}[A PWS three-term boundary certificate]
\label{r15allps:certificate}
For a PWS form $\omega$ of bidegree $(p,q)$, such a certificate
is a strict actual equality
\[
                 \omega=d'\mu+d''\nu+D\psi
\tag{PA15.1}\label{r15allps:three-term}
\]
with global PWS inputs of degrees
$(p-1,q)$, $(p,q-1)$, and $(p-1,q-1)$, respectively.
Negative degrees mean zero. Every such form is an actual Aeppli
boundary. The individual derivatives in
\eqref{r15allps:three-term} need not be PWS; only their combined
sum is required to be the given PWS form.
\end{definition}

\begin{lemma}[Finite descent of this certificate]
\label{r15allps:finite}
Let $q:V\to W$ be finite dominant between normal integral
projective $K$-varieties, of native degree $e$.
Let $\omega\in\operatorname{PS}^{p,q}(W)$.
If $q^*\omega$ has a PWS three-term boundary certificate
on $V$, then $\omega$ has one on $W$.
\end{lemma}

\begin{proof}
The frozen all-bidegree finite PWS trace is defined on
$q^*\omega$ and on all three supplied primitive inputs.
Its native projection identity gives
$\operatorname{Tr}^{\rm PS}_q(q^*\omega)=e\omega$.
Push the certificate as currents and commute proper current
differentiation. The prescribed all-actual-test trace pairings
and current injectivity give
\[
 \omega
 =d'\bigl(e^{-1}\operatorname{Tr}^{\rm PS}_q\mu\bigr)
  +d''\bigl(e^{-1}\operatorname{Tr}^{\rm PS}_q\nu\bigr)
  +D\bigl(e^{-1}\operatorname{Tr}^{\rm PS}_q\psi\bigr).
\tag{PA15.2}\label{r15allps:trace-certificate}
\]
The exact source is
\path{round14/foundations_audit/finite_pws_descent.tex},
Theorem \texttt{r14f:pws-trace} and its derivative-current
consequence. No trace of an arbitrary actual primitive or
of a separate positive-delta-codimension derivative is used.
\end{proof}

Choose a smooth curvature $h$ of $\mathcal O_{\mathbf P^1}(1)$.
It is separately closed, satisfies $h^2=0$ by dimension, and
has native mass one over $K$.

\begin{lemma}[A constructive decomposition on a product of projective lines]
\label{r15allps:product}
Let $R_m=(\mathbf P_K^1)^m$ and write $h_i$ for the pullback
of $h$ from its $i$th factor. Put $h_I=\prod_{i\in I}h_i$,
with $h_\varnothing=1$. Every $D$-closed
$\alpha\in\operatorname{PS}^{p,q}(R_m)$ has the following
strict actual decomposition, with a PWS three-term boundary
certificate on the right:
\[
 \alpha=
 \begin{cases}
 \displaystyle\sum_{|I|=k}c_I h_I+d'\mu+d''\nu+D\psi,
                                      &p=q=k,\quad0\leq k\leq m,\\
 d'\mu+d''\nu+D\psi,                    &p\ne q.
 \end{cases}
\tag{PA15.3}\label{r15allps:product-decomposition}
\]
Degrees above dimension are zero. The coefficients in the
diagonal case are unique even modulo all actual Aeppli
boundaries, and are given by
\[
                         c_I=\int_{R_m}\alpha h_{I^c}.
\tag{PA15.4}\label{r15allps:product-period}
\]
\end{lemma}

\begin{proof}
Consider $\pi:R_{m-1}\times\mathbf P^1\to R_{m-1}$ and
let $h_m$ be its last-factor curvature. Apply the proved
all-bidegree PWS incidence homotopy
Lemma~\ref{r15s:general-homotopy} to this trivial rank-two
bundle. Choose its quotient-line metric so that the incidence
curvature $\nu$ there is exactly $-h_m$.
Projective-line base change gives $C\alpha=\pi^*I\alpha$,
where
\[
              I\alpha=\pi_*\alpha,\qquad
              J\alpha=\pi_*(h_m\alpha).
\]
These are $D$-closed PWS forms of degrees $(p-1,q-1)$ and
$(p,q)$. The strict homotopy becomes
\[
 \alpha=\pi^*(J\alpha)+h_m\pi^*(I\alpha)
               +D\mathcal G\alpha+d''\mathcal P\alpha
                                      -d'\mathcal Q\alpha,
\tag{PA15.5}\label{r15allps:product-homotopy}
\]
with every transgression input PWS in its indicated degree.

Induct on $m$, starting with the point. Pull back the
certificates for $I\alpha$ and $J\alpha$. Multiplication by
the smooth separately closed $h_m$ preserves their three-term
form and their PWS inputs. Thus the induction yields
\eqref{r15allps:product-decomposition}. A scalar on the point
contributes only in bidegree $(0,0)$, so a monomial can occur
only when $p=q=k$ and it involves exactly $k$ distinct factors.

To check uniqueness in the full actual quotient, pair the
diagonal formula with $h_{I^c}$. The product $h_Jh_{I^c}$
is zero unless $J=I$, because otherwise one factor is squared.
For $J=I$ its native integral is one by iterated projective
fiber-degree calibration. Every actual Aeppli boundary, including
the $D\psi$ term, has zero pairing with this separately closed
test by actual Stokes. This proves
\eqref{r15allps:product-period} and the asserted uniqueness.
\end{proof}

\begin{definition}[PWS-represented groups in arbitrary bidegree]
\label{r15allps:group}
For a smooth integral projective variety $V/K$, let
\[
 H_{\mathrm A,\mathrm{PSrep}}^{p,q}(V)
   =\operatorname{im}\left(
       \ker(D:\operatorname{PS}^{p,q}(V)\to B^{p+1,q+1}(V))
        \longrightarrow H_{\mathrm A,\mathrm{act}}^{p,q}(V)
                         \right).
\tag{PA15.6}\label{r15allps:PSrep}
\]
This agrees with the earlier definitions in diagonal degrees
two and three; its quotient uses all actual boundary coefficients.
\end{definition}

\begin{theorem}[All PWS-represented groups of projective space]
\label{r15allps:projective-space}
For every $m\geq0$ and all nonnegative $p,q$,
\[
 H_{\mathrm A,\mathrm{PSrep}}^{p,q}(\mathbf P_K^m)
   =
   \begin{cases}
      \mathbf R[h^k],&p=q=k,\quad 0\leq k\leq m,\\
      0,&p\ne q\ \hbox{or}\ \max(p,q)>m.
   \end{cases}
\tag{PA15.7}\label{r15allps:Pm}
\]
Here $h$ may be any smooth curvature of $\mathcal O(1)$.
More precisely, every $D$-closed PWS representative of a
diagonal class differs from its scalar multiple of $h^k$
by a PWS three-term boundary certificate, and every off-diagonal
$D$-closed PWS form has such a certificate.
\end{theorem}

\begin{proof}
The cases $m=0$ and forms above dimension are immediate.
For $m\geq1$, use the finite symmetric-product morphism
\[
 q_m:(\mathbf P^1)^m\longrightarrow
               \operatorname{Sym}^m(\mathbf P^1)\simeq\mathbf P^m,
                         \qquad\deg_K q_m=m!.
\tag{PA15.8}\label{r15allps:symmetric-cover}
\]
It is the binary-form multiplication map, so
$q_m^*\mathcal O_{\mathbf P^m}(1)
 \simeq\bigotimes_{i=1}^m\operatorname{pr}_i^*
                  \mathcal O_{\mathbf P^1}(1)$.
The finite map and this native degree are the same ones
already used for PWS integration in
Section~\ref{r15ps:section}.

Let $\alpha$ be $D$-closed PWS on $\mathbf P^m$.
The pullback $q_m^*\alpha$ is PWS and is invariant under
all permutations of the factors. Apply
Lemma~\ref{r15allps:product}. In diagonal degree $k$, the
uniqueness of its actual Aeppli coefficients makes all
$c_I$ with $|I|=k$ equal to one real number $c$.
Indeed permutations act transitively on these subsets,
preserve the pullback form, and carry PWS certificates to
PWS certificates. In off-diagonal degree there are no
monomial terms.

Let $H=\sum_i h_i$ on the source. The ratio between the
pullback smooth metric defining $q_m^*h$ and the product
smooth metric defining $H$ is a global smooth scalar $u$,
with sign chosen so that $q_m^*h-H=Du$.
For $k\geq1$,
\[
 (q_m^*h)^k-H^k
   =D\left(u\sum_{j=0}^{k-1}
                    (q_m^*h)^jH^{k-1-j}\right),
       \qquad
 H^k=k!\sum_{|I|=k}h_I.
\tag{PA15.9}\label{r15allps:metric}
\]
The primitive inside the first $D$ is smooth, and every
curvature factor is separately closed. The case $k=0$
is the identity $1=1$. Therefore
$q_m^*(\alpha-(c/k!)h^k)$ has a PWS three-term boundary
certificate in diagonal degree, with the evident value
$c$ when $k=0$. In off-diagonal degree $q_m^*\alpha$
already has that certificate. Lemma~\ref{r15allps:finite}
descends it using the native degree $m!$ and the traces
of its three PWS primitive inputs. This proves spanning
and the stronger certificate statement.

Finally $[h^k]\ne0$: pair with the separately closed
$h^{m-k}$ and use $\int_{\mathbf P^m}h^m=1$.
Actual Aeppli boundaries have zero such pairing.
This proves the claimed dimension in the full actual
quotient, not only in a smaller PWS boundary quotient.
\end{proof}

\begin{corollary}[Finite constant fields and original tests]
\label{r15allps:constant-field}
For finite $L/K$, the same formula holds for
$\mathbf P_L^m$ regarded as a $K$-variety, in the original
$K$-actual theory. Its nonzero diagonal generator has native
top pairing $[L:K]$. On $(\mathbf P_L^1)^m$, the period
formula \eqref{r15allps:product-period} is divided by $[L:K]$.
\end{corollary}

\begin{proof}
Run the entire argument on these smooth integral projective
$K$-schemes. Their projective-line operators and finite traces
are the original-field ones. Each relative projective-line
mass is one, while the final zero-dimensional native mass is
$[L:K]$. The symmetric-product map between the two varieties
with constant field $L$ still has native degree $m!$.
No equality of test categories over different fields is inferred.
\end{proof}

\paragraph{Full actual groups and the later degree-two result.}
The theorem computes the image of $D$-closed PWS forms in
the actual Aeppli quotient. It does not prove
$H_{\mathrm A,\mathrm{act}}^{p,q}(\mathbf P^m)
 =H_{\mathrm A,\mathrm{PSrep}}^{p,q}(\mathbf P^m)$
when a new positive-delta-codimension representative problem
arises. The full actual groups in degrees $(1,1)$ and
$(2,2)$ are proved by the separate incidence and
fixed-section arguments; see Corollary~\ref{r15a2:full-HA22}
for the latter. The full actual diagonal groups in degrees
$(k,k)$ with $k\geq3$ remain open. The retained separate
$D$-primitive in the certificate is essential; its first
derivative need not be PWS. The theorem does not give a trace
on arbitrary actual forms, replace every actual nonclosed
primitive by PWS, or prove the general strong conjecture G.
